\section{Introduction}

Neural networks trained with standard objectives repeatedly exhibit behaviors associated with probabilistic inference: soft clustering, prototype specialization, uncertainty tracking, and mixture-model dynamics. These phenomena appear across architectures---in attention mechanisms, classification heads, and energy-based models---yet their origin has remained unclear. Are they emergent properties of scale? Architectural accidents? Or something more fundamental? This paper argues that such behaviors are none of these. They are necessary consequences of the geometry of common objective functions.

\subsection{The Puzzle}

Consider the range of phenomena that appear, unbidden, in trained neural networks. Attention heads in transformers learn to specialize, each routing information for distinct semantic roles. Classification networks partition representation space into regions that behave like mixture components. Deep networks trained on noisy data exhibit robustness patterns reminiscent of Bayesian inference, down-weighting outliers and tracking uncertainty across inputs. These behaviors emerge without explicit probabilistic modeling, without mixture-model architectures, and without any algorithm resembling expectation-maximization.

The standard explanations are unsatisfying. One view holds that these are emergent properties of scale---that sufficient parameters and data somehow give rise to statistical structure. Another treats them as architectural coincidences, artifacts of specific design choices like softmax normalization or residual connections. A third offers loose analogies: attention is ``like'' soft clustering; cross-entropy ``resembles'' a mixture model. None of these explanations answers the deeper question: \emph{why do these specific behaviors appear, rather than others?} And why do they appear so reliably, across such different architectures and tasks?

\subsection{Recent Evidence}

Recent work by \citet{aggarwal2025geometry,aggarwal2025gradient} has sharpened this puzzle considerably. In controlled experimental settings---``Bayesian wind tunnels'' where the true posterior is analytically known---small transformers match the analytic posterior's entropy, position by position, to within $10^{-3}$--$10^{-4}$ bits: agreement at the resolution of the arithmetic. Capacity-matched MLPs trained under identical conditions fail catastrophically, suggesting that the phenomenon depends on the inductive biases of attention rather than optimization alone.

More striking still is what Agarwal et al.\ find in the gradient dynamics. Attention weights stabilize early in training while value vectors continue to refine---a two-timescale structure that mirrors the E-step and M-step of classical expectation-maximization. Values receive updates weighted by attention, precisely as prototypes receive updates weighted by responsibilities in mixture models. The authors provide a complete first-order analysis showing that this structure is not incidental but systematic.

A scoping distinction, which the authors make explicit in later revisions of both papers, is essential to everything that follows. Two different processes are in play. The \emph{Bayesian inference} their wind tunnels demonstrate is in-context computation: a frozen, trained network tracking the posterior predictive over latent task variables as a sequence unfolds, at inference time. The \emph{EM-like dynamics} are a training-time phenomenon: their gradient analysis, in their words, ``lives at the EM/SGD level,'' explaining how cross-entropy training sculpts the geometry that the frozen network later uses. The present paper's claims live entirely at the second level---training time. We analyze what gradient descent does, not what the learned function computes in context.

Within that training-time level, Agarwal et al.\ are careful about the status of the EM correspondence. Their gradient laws are derived, not analogized---the mechanics are exact consequences of cross-entropy. What they decline is the probabilistic interpretation: because value updates are driven by the backpropagated error signal rather than by observed data, values move to explain the error geometry rather than to maximize a likelihood over inputs, and so, in their words, ``the analogy is structural rather than variational.'' What their account leaves open is the probabilistic status of the training dynamics: whether the responsibilities that gate learning are posteriors of anything.

\subsection{This Paper}

This paper supplies that status, and delimits it. We show that for objectives with log-sum-exp structure over distances or energies, the gradient of the loss with respect to each distance is exactly the posterior responsibility of the corresponding component. This is not an approximation, not a resemblance, not an analogy. It is an algebraic identity:

\begin{equation}
\frac{\partial L}{\partial d_j} = -r_j
\end{equation}

The identity itself is classical---it is Fisher's identity from the EM literature, specialized to log-sum-exp objectives---and the precise relationship between gradient descent and EM has been characterized before \citep{xu1996convergence,neal1998view}. What has not been drawn out is where this identity lives: standard neural objectives instantiate it, unmodified, over the distance-based representations that ordinary networks compute. When the log-sum-exp is the loss, the EM correspondence is therefore not structural but variational: the responsibilities are posterior probabilities of an implicit mixture, and gradient descent performs expectation-maximization implicitly, continuously rather than in discrete alternating steps. The forward pass computes unnormalized likelihoods; normalization yields responsibilities; backpropagation delivers responsibility-weighted updates to parameters. No auxiliary latent variables are introduced. No inference algorithm is invoked. The gradient \emph{is} the responsibility. When the softmax instead sits inside the network---attention---the same calculus recovers exactly the structural laws of \citet{aggarwal2025gradient}, and explains why the variational reading applies at the loss level and becomes structural in the interior (Section 4.2).

This reframes the relationship between optimization and inference. Inference is not a separate algorithmic layer added on top of learning. It is not a post-hoc interpretation of learned representations. Under the objectives we analyze, inference and optimization are the same computation viewed at different levels of abstraction. The Bayesian structure observed by Agarwal et al.\ is not an emergent property that happens to appear; it is forced by the geometry of the loss.

\subsection{Contributions}

The contributions of this paper are deliberately narrow and can be summarized as three claims at increasing levels of abstraction.

\textbf{One theorem.} For any objective of the form $L = \log \sum_j \exp(-d_j)$, the gradient with respect to each distance satisfies $\partial L / \partial d_j = -r_j$, where $r_j$ is the normalized exponential---the responsibility of component $j$. The identity is a specialization of Fisher's identity from the classical EM literature (Section 5.2); it requires no conditions beyond differentiability. What is new is not the identity but its address: standard neural objectives---cross-entropy, attention, mixture likelihoods---instantiate it without modification.

\textbf{One interpretation.} This identity implies that gradient descent on distance-based log-sum-exp objectives performs implicit expectation-maximization. The E-step is the forward pass; the M-step is the parameter update; responsibilities are never explicitly computed because they \emph{are} the gradients. EM is not approximated by neural training---it is realized by it.

\textbf{One unification.} The same mechanism manifests across three regimes depending on the constraints imposed. In the \emph{unsupervised} regime, responsibilities are fully latent and prototypes compete freely. In the \emph{conditional} regime---attention---responsibilities are recomputed per query over a shared prototype family. In the \emph{constrained} regime---cross-entropy classification---supervision clamps one responsibility while competition persists among alternatives. These are not three different phenomena. They are one phenomenon under different boundary conditions.
